% ---------------------------------------------------------------------
% Das Dokument kompiliert mit pdflatex und ist auf Basis
% von Koma-Script entstanden.
%
% Autor des Templates (für Anmerkungen):
% Michael von Riegen, riegen@informatik.uni-hamburg.de
%
% Einzelne Code-Teile für das Titelblatt sind aus dem Template
% von Benjamin Kirchheim entnommen.
%
% 25.05.09, Frank Langanke: Vorlage auf aktuelle KOMA-Version aktualisiert
% 26.05.09, Michael von Riegen: Anmerkung --> aktuelles Koma-Script ist nötig!
% 17.10.2016 neues Uni logo
% ---------------------------------------------------------------------
\documentclass[11pt,DIV=15,BCOR=20mm,bibliography=totoc]{scrbook}

% Import von Paketen und Optionen die das gesamte Dokument betreffen
% sind in myPreamble.sty ausgelagert.
\usepackage{myPreamble}

% Arbeitet man nur an einem Kapitel, wird durch folgenden Befehl nur dieses eingebunden.
% Spart manuelles auskommentieren von vielen include-Befehlen;
% hat keine Auswirkung auf input-Befehle
% \includeonly{kapitel1}

\begin{document}

% TITELSEITE
\begin{titlepage}

	% Fehler "destination with the same identifier" unterdrücken...
  \setcounter{page}{-1}

	% Titelseite
	\begin{figure}[h]
		\begin{minipage}[b]{62mm}
			\includegraphics[width=62mm]{images/unilogo}
		\end{minipage}
		\hspace{4cm}
		\begin{minipage}[b]{59mm}
			\includegraphics[width=59mm]{images/minlogo}
		\end{minipage}
	\end{figure}

	\vfill

	\begin{center}
		% Diplomarbeit
		\noindent { \Large
			Bachelor's Thesis \\
		}
		\vspace{14mm}
		% Titel
		\noindent \textbf{\huge
		A Recommender Framework for Skills Management \\
		%   Make Skills Management Great Again\\
		}
		\vspace{14mm}
    \noindent { \Large
			In Cooperation With SinnerSchrader \\
			% Design and Implementation of a Skills Management Web Application \\
		}
	\end{center}

	\vfill

	\noindent \textbf{Torben Reetz} \\
	\noindent \rule{\textwidth}{0.4mm}
	\noindent{\textrm{3reetz@informatik.uni-hamburg.de}} \\
	\noindent{\textrm{B.Sc. Software Systems Development}} \\
	\noindent{\textrm{Matr.-Nr. 6524064}} \\
	\noindent{\textrm{Semester 07}} \\
	\begin{tabbing}
	\hspace{8em} \=  \kill
	Primary Supervisor: Prof. Dr. Walid Maalej \\
	Secondary Supervisor: Prof. Dr. Eva Bittner \\
	~ \\
	Date: xx.xx.2017
	\end{tabbing}

	% Rückseite der Titelseite mit Zitat
	% \newpage
	% \thispagestyle{empty}
	% \setcounter{page}{0}

	% wenn man Lust auf ein Zitat hat...
	% ... ansonsten auskommentieren
	% ~\\ \vfill \noindent
	% A distributed system is one where the failure of some \\
	% computer I've never heard of can keep me from getting my work done. \\
	% \textit{-- Leslie Lamport}
\end{titlepage}


% VERZEICHNISSE (Inhaltsverzeichnis, Abkürzungen)
% Vorspann einleiten --> Seitennummerierung römisch
\frontmatter

% Contentseses
\tableofcontents
\cleardoublepage

% Hauptteil einleiten --> Seitennummerierung wieder arabisch
\mainmatter

% -----------------------------------------------------------------------
% -----------------------------------------------------------------------
% -----------------------------------------------------------------------
% Einleitung
% -----------------------------------------------------------------------
% -----------------------------------------------------------------------
% -----------------------------------------------------------------------
\chapter{Einleitung}

Hier kommt eine Quellenangabe \cite{Gray1981}, \cite{Cerami2002}. Weiterhin sieht man hier einen Link auf Kapitel \ref{ch:1}. Lorem ipsum dolor sit amet, consectetuer\footnote{Hier ist eine Fußnote!} adipiscing elit. Integer quis lectus eget purus auctor sollicitudin. Sed tempus, leo quis nonummy iaculis, tellus justo bibendum mi, eu facilisis est sapien non eros. Ut tincidunt. Proin eleifend tristique est. Nulla facilisi. Cras ante mauris, facilisis at, fringilla quis, venenatis ornare, libero. Nam viverra varius nibh. Suspendisse potenti. Sed velit quam, euismod quis, sagittis eu, sodales sit amet, dui. Sed sed diam. Vivamus tincidunt quam a eros. Ut quis metus. Lorem ipsum dolor sit amet, consectetuer adipiscing elit. Cum sociis natoque penatibus et magnis dis parturient montes, nascetur ridiculus mus. Sed ut tortor. Vestibulum ante ipsum primis in faucibus orci luctus et ultrices posuere cubilia Curae; Curabitur vulputate nibh a tortor. Nullam scelerisque risus nonummy urna tempor sagittis.

\noindent Hier kommt ein Listing \ref{lst:soap}.

\begin{center}
\begin{lstlisting}[caption={SOAP Anfrage an einen HalloWelt-Web-Service},label=lst:soap,language=XML,label={lst:soap}]
<?xml version='1.0' encoding='UTF-8'>
<SOAP-ENV:Envelope (*@\label{lst:soapEnv}@*)
  xmlns:SOAP-ENV="http://schemas.xmlsoap.org/soap/envelope/"
  xmlns:xsi="http://www.w3.org/2001/XMLSchema-instance"
  xmlns:xsd="http://www.w3.org/2001/XMLSchema"
  xmlns:ns1="http://localhost/wsdl/HalloWeltService.wsdl">
  
  <SOAP-ENV:Body>(*@\label{lst:soapBody}@*)
  	<ns1:gruss>
  		<name xsi:type="xsd:string">
  			Michael
  		</name>
  	</ns1:gruss>
  </SOAP-ENV:Body>

</SOAP-ENV:Envelope>
\end{lstlisting}
\end{center}

Donec tincidunt. Cras tempus purus et quam. Duis semper. Vestibulum eu ligula. Ut non odio. Ut posuere ligula fringilla neque. Cras sodales justo eu massa. Nullam mauris nulla, fermentum et, fringilla sed, imperdiet vitae, turpis. In vehicula fermentum neque. Cum sociis natoque penatibus et magnis dis parturient montes, nascetur ridiculus mus. Quisque orci. Nunc odio. Nam consectetuer egestas augue.

\noindent Hier kommt eine etwas komplizierte Tabelle:

\begin{table}[h]
	\caption[Transaktionale Ablaufmodelle im Vergleich]{Transaktionale Ablaufmodelle im Vergleich; teilweise aus \cite{DBLP:books/infix/Schwarz99}; das Wort "`Transaktion"' wird in der Tabelle als TA abgekürzt}
	\label{tab:tamodelleVergleich}
	\centering
		\begin{tabular*}{\textwidth}{| l@{\extracolsep\fill} r || c | c | c || c | c || c || c || c | c | c || c | c |}
			\hline
			
			Modelle & 
				\rotatebox{90}{Eigenschaften} &
					\rotatebox{90}{vital} &
						\rotatebox{90}{non-vital} &
							\rotatebox{90}{Alternativtransaktion } &				
								\rotatebox{90}{sequenziell} &
									\rotatebox{90}{parallel} &
  									\rotatebox{90}{unabhängig} &
	    								\rotatebox{90}{temporal} &
			    							\rotatebox{90}{geschlossen geschachtelt } &
					    						\rotatebox{90}{offen geschachtelt} &
							    					\rotatebox{90}{Kompensation} &
									    				\rotatebox{90}{endliche Schachtelung } &
											    			\rotatebox{90}{unendliche Schachtelung } \\
			
			\hline
			\hline
			
			\multicolumn{2}{|l||}{Geschl. gesch. TA} &
				\checkmark &
				 (\checkmark) &
				 	(\checkmark) &
				 		\checkmark &
				 			\checkmark &
				 			
				 			  - &
				 			  - &
				 			
				 				\checkmark &
				 				  - &
				 				  	-	&
				 				  		\checkmark &
				 				  			\checkmark \\
			\hline

			\multicolumn{2}{|l||}{Offen gesch. TA} &
				\checkmark &
				 (\checkmark) &
				 	(\checkmark) &
				 		\checkmark &
				 			\checkmark &
				 			 
				 			  - &
				 			  - &
				 			
				 				- &
				 				  \checkmark &
				 				  	\checkmark &
				 				  		\checkmark &
				 				  			\checkmark \\
			\hline

			\multicolumn{2}{|l||}{Sagas} &
				\checkmark &
				 - &
				 	- &
				 		\checkmark &
				 			- &
				 			
				 			  - &
				 			  - &
				 			
				 				- &
				 				  \checkmark &
				 				  	\checkmark &	
				 				  		\checkmark &
				 				  			- \\
			\hline

			\multicolumn{2}{|l||}{Flex-TA} &
				\checkmark &
				 \checkmark &
				 	\checkmark &
				 		\checkmark &
				 			\checkmark &
				 			
				 			- &
				 			\checkmark &
				 			 
				 			
				 				\checkmark &
				 				  \checkmark &
				 				  	\checkmark &
				 				  		\checkmark &
				 				  			- \\
			
			\hline

			\multicolumn{2}{|l||}{ConTracts} &
				\checkmark &
				 \checkmark &
				 	- &
				 		\checkmark &
				 			\checkmark &
				 			
				 			 -&
				 			 -&
				 			
				 				\checkmark &
				 				  \checkmark &
				 				  	\checkmark &
				 				  		\checkmark &
				 				  			\checkmark \\

			\hline

			\multicolumn{2}{|l||}{Business-TA} &
				\checkmark &
				 \checkmark &
				 	\checkmark &
				 		\checkmark &
				 			\checkmark &

  			 			 \checkmark &
				 			 \checkmark &
				 			  
				 				\checkmark &
				 				  \checkmark &
				 				  	\checkmark &
				 				  		\checkmark &
				 				  			\checkmark \\			 				  			
			\hline
				
		\end{tabular*}
\end{table}

Vestibulum sed odio. Mauris semper placerat ipsum. Nulla facilisi. In rutrum, tortor in rhoncus venenatis, diam enim lacinia urna, vel egestas nunc sapien ut nisl. Fusce placerat posuere est. Maecenas augue. In gravida leo vel massa. Aliquam eu nisi. Cras tincidunt aliquam arcu. Pellentesque habitant morbi tristique senectus et netus et malesuada fames ac turpis egestas. Aenean quis libero condimentum libero porttitor vestibulum. Pellentesque a nulla nec ipsum tristique molestie.

Curabitur sagittis quam vitae tortor. Nulla id pede. Praesent rutrum, ipsum quis porttitor pretium, nisl ipsum gravida est, commodo pulvinar augue nibh ac felis. Aliquam erat volutpat. Aenean ullamcorper feugiat dolor. Donec nulla mauris, elementum sed, varius ut, blandit vitae, purus. Aenean in nunc in justo sodales fringilla. Proin pulvinar. Maecenas ullamcorper. Vivamus dapibus consectetuer arcu. Nam dictum, erat non adipiscing pretium, lorem tellus pulvinar dui, ut condimentum massa est sit amet urna. Nulla facilisi.

Vivamus in dui. Morbi tristique nibh in purus. Nam erat. Curabitur euismod lacinia purus. Maecenas tempor libero quis dolor. Quisque gravida lorem vitae libero. Aliquam sodales, risus in tempor imperdiet, massa nisi auctor massa, vitae laoreet nunc erat in sapien. Etiam eu sapien nec enim consectetuer accumsan. Nunc non quam. Cras placerat justo quis nibh nonummy sollicitudin. Aenean vitae felis. Mauris id mi nec ligula ultricies vestibulum. Suspendisse luctus pretium justo. 

\section{Motivation}
\label{ch:1}
Lorem ipsum dolor sit amet, consectetuer adipiscing elit. Integer quis lectus eget purus auctor sollicitudin. Sed tempus, leo quis nonummy iaculis, tellus justo bibendum mi, eu facilisis est sapien non eros. Ut tincidunt. Proin eleifend tristique est. Nulla facilisi. Cras ante mauris, facilisis at, fringilla quis, venenatis ornare, libero. Nam viverra varius nibh. Suspendisse potenti. Sed velit quam, euismod quis, sagittis eu, sodales sit amet, dui. Sed sed diam. Vivamus tincidunt quam a eros. Ut quis metus. Lorem ipsum dolor sit amet, consectetuer adipiscing elit. Cum sociis natoque penatibus et magnis dis parturient montes, nascetur ridiculus mus. Sed ut tortor. Vestibulum ante ipsum primis in faucibus orci luctus et ultrices posuere cubilia Curae; Curabitur vulputate nibh a tortor. Nullam scelerisque risus nonummy urna tempor sagittis.

Donec tincidunt. Cras tempus purus et quam. Duis semper. Vestibulum eu ligula. Ut non odio. Ut posuere ligula fringilla neque. Cras sodales justo eu massa. Nullam mauris nulla, fermentum et, fringilla sed, imperdiet vitae, turpis. In vehicula fermentum neque. Cum sociis natoque penatibus et magnis dis parturient montes, nascetur ridiculus mus. Quisque orci. Nunc odio. Nam consectetuer egestas augue.

Vestibulum sed odio. Mauris semper placerat ipsum. Nulla facilisi. In rutrum, tortor in rhoncus venenatis, diam enim lacinia urna, vel egestas nunc sapien ut nisl. Fusce placerat posuere est. Maecenas augue. In gravida leo vel massa. Aliquam eu nisi. Cras tincidunt aliquam arcu. Pellentesque habitant morbi tristique senectus et netus et malesuada fames ac turpis egestas. Aenean quis libero condimentum libero porttitor vestibulum. Pellentesque a nulla nec ipsum tristique molestie.

Curabitur sagittis quam vitae tortor. Nulla id pede. Praesent rutrum, ipsum quis porttitor pretium, nisl ipsum gravida est, commodo pulvinar augue nibh ac felis. Aliquam erat volutpat. Aenean ullamcorper feugiat dolor. Donec nulla mauris, elementum sed, varius ut, blandit vitae, purus. Aenean in nunc in justo sodales fringilla. Proin pulvinar. Maecenas ullamcorper. Vivamus dapibus consectetuer arcu. Nam dictum, erat non adipiscing pretium, lorem tellus pulvinar dui, ut condimentum massa est sit amet urna. Nulla facilisi.

Vivamus in dui. Morbi tristique nibh in purus. Nam erat. Curabitur euismod lacinia purus. Maecenas tempor libero quis dolor. Quisque gravida lorem vitae libero. Aliquam sodales, risus in tempor imperdiet, massa nisi auctor massa, vitae laoreet nunc erat in sapien. Etiam eu sapien nec enim consectetuer accumsan. Nunc non quam. Cras placerat justo quis nibh nonummy sollicitudin. Aenean vitae felis. Mauris id mi nec ligula ultricies vestibulum. Suspendisse luctus pretium justo. 

Vestibulum sed odio. Mauris semper placerat ipsum. Nulla facilisi. In rutrum, tortor in rhoncus venenatis, diam enim lacinia urna, vel egestas nunc sapien ut nisl. Fusce placerat posuere est. Maecenas augue. In gravida leo vel massa. Aliquam eu nisi. Cras tincidunt aliquam arcu. Pellentesque habitant morbi tristique senectus et netus et malesuada fames ac turpis egestas. Aenean quis libero condimentum libero porttitor vestibulum. Pellentesque a nulla nec ipsum tristique molestie.

Curabitur sagittis quam vitae tortor. Nulla id pede. Praesent rutrum, ipsum quis porttitor pretium, nisl ipsum gravida est, commodo pulvinar augue nibh ac felis. Aliquam erat volutpat. Aenean ullamcorper feugiat dolor. Donec nulla mauris, elementum sed, varius ut, blandit vitae, purus. Aenean in nunc in justo sodales fringilla. Proin pulvinar. Maecenas ullamcorper. Vivamus dapibus consectetuer arcu. Nam dictum, erat non adipiscing pretium, lorem tellus pulvinar dui, ut condimentum massa est sit amet urna. Nulla facilisi.

Vivamus in dui. Morbi tristique nibh in purus. Nam erat. Curabitur euismod lacinia purus. Maecenas tempor libero quis dolor. Quisque gravida lorem vitae libero. Aliquam sodales, risus in tempor imperdiet, massa nisi auctor massa, vitae laoreet nunc erat in sapien. Etiam eu sapien nec enim consectetuer accumsan. Nunc non quam. Cras placerat justo quis nibh nonummy sollicitudin. Aenean vitae felis. Mauris id mi nec ligula ultricies vestibulum. Suspendisse luctus pretium justo. 

Vestibulum sed odio. Mauris semper placerat ipsum. Nulla facilisi. In rutrum, tortor in rhoncus venenatis, diam enim lacinia urna, vel egestas nunc sapien ut nisl. Fusce placerat posuere est. Maecenas augue. In gravida leo vel massa. Aliquam eu nisi. Cras tincidunt aliquam arcu. Pellentesque habitant morbi tristique senectus et netus et malesuada fames ac turpis egestas. Aenean quis libero condimentum libero porttitor vestibulum. Pellentesque a nulla nec ipsum tristique molestie.

Curabitur sagittis quam vitae tortor. Nulla id pede. Praesent rutrum, ipsum quis porttitor pretium, nisl ipsum gravida est, commodo pulvinar augue nibh ac felis. Aliquam erat volutpat. Aenean ullamcorper feugiat dolor. Donec nulla mauris, elementum sed, varius ut, blandit vitae, purus. Aenean in nunc in justo sodales fringilla. Proin pulvinar. Maecenas ullamcorper. Vivamus dapibus consectetuer arcu. Nam dictum, erat non adipiscing pretium, lorem tellus pulvinar dui, ut condimentum massa est sit amet urna. Nulla facilisi.

Vivamus in dui. Morbi tristique nibh in purus. Nam erat. Curabitur euismod lacinia purus. Maecenas tempor libero quis dolor. Quisque gravida lorem vitae libero. Aliquam sodales, risus in tempor imperdiet, massa nisi auctor massa, vitae laoreet nunc erat in sapien. Etiam eu sapien nec enim consectetuer accumsan. Nunc non quam. Cras placerat justo quis nibh nonummy sollicitudin. Aenean vitae felis. Mauris id mi nec ligula ultricies vestibulum. Suspendisse luctus pretium justo. 

Vestibulum sed odio. Mauris semper placerat ipsum. Nulla facilisi. In rutrum, tortor in rhoncus venenatis, diam enim lacinia urna, vel egestas nunc sapien ut nisl. Fusce placerat posuere est. Maecenas augue. In gravida leo vel massa. Aliquam eu nisi. Cras tincidunt aliquam arcu. Pellentesque habitant morbi tristique senectus et netus et malesuada fames ac turpis egestas. Aenean quis libero condimentum libero porttitor vestibulum. Pellentesque a nulla nec ipsum tristique molestie.

Curabitur sagittis quam vitae tortor. Nulla id pede. Praesent rutrum, ipsum quis porttitor pretium, nisl ipsum gravida est, commodo pulvinar augue nibh ac felis. Aliquam erat volutpat. Aenean ullamcorper feugiat dolor. Donec nulla mauris, elementum sed, varius ut, blandit vitae, purus. Aenean in nunc in justo sodales fringilla. Proin pulvinar. Maecenas ullamcorper. Vivamus dapibus consectetuer arcu. Nam dictum, erat non adipiscing pretium, lorem tellus pulvinar dui, ut condimentum massa est sit amet urna. Nulla facilisi.

Vivamus in dui. Morbi tristique nibh in purus. Nam erat. Curabitur euismod lacinia purus. Maecenas tempor libero quis dolor. Quisque gravida lorem vitae libero. Aliquam sodales, risus in tempor imperdiet, massa nisi auctor massa, vitae laoreet nunc erat in sapien. Etiam eu sapien nec enim consectetuer accumsan. Nunc non quam. Cras placerat justo quis nibh nonummy sollicitudin. Aenean vitae felis. Mauris id mi nec ligula ultricies vestibulum. Suspendisse luctus pretium justo. 

\section{Zielsetzung}
Lorem ipsum dolor sit amet, consectetuer adipiscing elit. Integer quis lectus eget purus auctor sollicitudin. Sed tempus, leo quis nonummy iaculis, tellus justo bibendum mi, eu facilisis est sapien non eros. Ut tincidunt. Proin eleifend tristique est. Nulla facilisi. Cras ante mauris, facilisis at, fringilla quis, venenatis ornare, libero. Nam viverra varius nibh. Suspendisse potenti. Sed velit quam, euismod quis, sagittis eu, sodales sit amet, dui. Sed sed diam. Vivamus tincidunt quam a eros. Ut quis metus. Lorem ipsum dolor sit amet, consectetuer adipiscing elit. Cum sociis natoque penatibus et magnis dis parturient montes, nascetur ridiculus mus. Sed ut tortor. Vestibulum ante ipsum primis in faucibus orci luctus et ultrices posuere cubilia Curae; Curabitur vulputate nibh a tortor. Nullam scelerisque risus nonummy urna tempor sagittis.

Donec tincidunt. Cras tempus purus et quam. Duis semper. Vestibulum eu ligula. Ut non odio. Ut posuere ligula fringilla neque. Cras sodales justo eu massa. Nullam mauris nulla, fermentum et, fringilla sed, imperdiet vitae, turpis. In vehicula fermentum neque. Cum sociis natoque penatibus et magnis dis parturient montes, nascetur ridiculus mus. Quisque orci. Nunc odio. Nam consectetuer egestas augue.

Vestibulum sed odio. Mauris semper placerat ipsum. Nulla facilisi. In rutrum, tortor in rhoncus venenatis, diam enim lacinia urna, vel egestas nunc sapien ut nisl. Fusce placerat posuere est. Maecenas augue. In gravida leo vel massa. Aliquam eu nisi. Cras tincidunt aliquam arcu. Pellentesque habitant morbi tristique senectus et netus et malesuada fames ac turpis egestas. Aenean quis libero condimentum libero porttitor vestibulum. Pellentesque a nulla nec ipsum tristique molestie.

Curabitur sagittis quam vitae tortor. Nulla id pede. Praesent rutrum, ipsum quis porttitor pretium, nisl ipsum gravida est, commodo pulvinar augue nibh ac felis. Aliquam erat volutpat. Aenean ullamcorper feugiat dolor. Donec nulla mauris, elementum sed, varius ut, blandit vitae, purus. Aenean in nunc in justo sodales fringilla. Proin pulvinar. Maecenas ullamcorper. Vivamus dapibus consectetuer arcu. Nam dictum, erat non adipiscing pretium, lorem tellus pulvinar dui, ut condimentum massa est sit amet urna. Nulla facilisi.

Vivamus in dui. Morbi tristique nibh in purus. Nam erat. Curabitur euismod lacinia purus. Maecenas tempor libero quis dolor. Quisque gravida lorem vitae libero. Aliquam sodales, risus in tempor imperdiet, massa nisi auctor massa, vitae laoreet nunc erat in sapien. Etiam eu sapien nec enim consectetuer accumsan. Nunc non quam. Cras placerat justo quis nibh nonummy sollicitudin. Aenean vitae felis. Mauris id mi nec ligula ultricies vestibulum. Suspendisse luctus pretium justo. 

\section{Eingrenzung des Themas}
\label{ch:eingrenzungThema}
Lorem ipsum dolor sit amet, consectetuer adipiscing elit. Integer quis lectus eget purus auctor sollicitudin. Sed tempus, leo quis nonummy iaculis, tellus justo bibendum mi, eu facilisis est sapien non eros. Ut tincidunt. Proin eleifend tristique est. Nulla facilisi. Cras ante mauris, facilisis at, fringilla quis, venenatis ornare, libero. Nam viverra varius nibh. Suspendisse potenti. Sed velit quam, euismod quis, sagittis eu, sodales sit amet, dui. Sed sed diam. Vivamus tincidunt quam a eros. Ut quis metus. Lorem ipsum dolor sit amet, consectetuer adipiscing elit. Cum sociis natoque penatibus et magnis dis parturient montes, nascetur ridiculus mus. Sed ut tortor. Vestibulum ante ipsum primis in faucibus orci luctus et ultrices posuere cubilia Curae; Curabitur vulputate nibh a tortor. Nullam scelerisque risus nonummy urna tempor sagittis.

Donec tincidunt. Cras tempus purus et quam. Duis semper. Vestibulum eu ligula. Ut non odio. Ut posuere ligula fringilla neque. Cras sodales justo eu massa. Nullam mauris nulla, fermentum et, fringilla sed, imperdiet vitae, turpis. In vehicula fermentum neque. Cum sociis natoque penatibus et magnis dis parturient montes, nascetur ridiculus mus. Quisque orci. Nunc odio. Nam consectetuer egestas augue.

Vestibulum sed odio. Mauris semper placerat ipsum. Nulla facilisi. In rutrum, tortor in rhoncus venenatis, diam enim lacinia urna, vel egestas nunc sapien ut nisl. Fusce placerat posuere est. Maecenas augue. In gravida leo vel massa. Aliquam eu nisi. Cras tincidunt aliquam arcu. Pellentesque habitant morbi tristique senectus et netus et malesuada fames ac turpis egestas. Aenean quis libero condimentum libero porttitor vestibulum. Pellentesque a nulla nec ipsum tristique molestie.

Curabitur sagittis quam vitae tortor. Nulla id pede. Praesent rutrum, ipsum quis porttitor pretium, nisl ipsum gravida est, commodo pulvinar augue nibh ac felis. Aliquam erat volutpat. Aenean ullamcorper feugiat dolor. Donec nulla mauris, elementum sed, varius ut, blandit vitae, purus. Aenean in nunc in justo sodales fringilla. Proin pulvinar. Maecenas ullamcorper. Vivamus dapibus consectetuer arcu. Nam dictum, erat non adipiscing pretium, lorem tellus pulvinar dui, ut condimentum massa est sit amet urna. Nulla facilisi.

Vivamus in dui. Morbi tristique nibh in purus. Nam erat. Curabitur euismod lacinia purus. Maecenas tempor libero quis dolor. Quisque gravida lorem vitae libero. Aliquam sodales, risus in tempor imperdiet, massa nisi auctor massa, vitae laoreet nunc erat in sapien. Etiam eu sapien nec enim consectetuer accumsan. Nunc non quam. Cras placerat justo quis nibh nonummy sollicitudin. Aenean vitae felis. Mauris id mi nec ligula ultricies vestibulum. Suspendisse luctus pretium justo. 

\section{Aufbau der Arbeit} 
Lorem ipsum dolor sit amet, consectetuer adipiscing elit. Integer quis lectus eget purus auctor sollicitudin. Sed tempus, leo quis nonummy iaculis, tellus justo bibendum mi, eu facilisis est sapien non eros. Ut tincidunt. Proin eleifend tristique est. Nulla facilisi. Cras ante mauris, facilisis at, fringilla quis, venenatis ornare, libero. Nam viverra varius nibh. Suspendisse potenti. Sed velit quam, euismod quis, sagittis eu, sodales sit amet, dui. Sed sed diam. Vivamus tincidunt quam a eros. Ut quis metus. Lorem ipsum dolor sit amet, consectetuer adipiscing elit. Cum sociis natoque penatibus et magnis dis parturient montes, nascetur ridiculus mus. Sed ut tortor. Vestibulum ante ipsum primis in faucibus orci luctus et ultrices posuere cubilia Curae; Curabitur vulputate nibh a tortor. Nullam scelerisque risus nonummy urna tempor sagittis.

Donec tincidunt. Cras tempus purus et quam. Duis semper. Vestibulum eu ligula. Ut non odio. Ut posuere ligula fringilla neque. Cras sodales justo eu massa. Nullam mauris nulla, fermentum et, fringilla sed, imperdiet vitae, turpis. In vehicula fermentum neque. Cum sociis natoque penatibus et magnis dis parturient montes, nascetur ridiculus mus. Quisque orci. Nunc odio. Nam consectetuer egestas augue.

Vestibulum sed odio. Mauris semper placerat ipsum. Nulla facilisi. In rutrum, tortor in rhoncus venenatis, diam enim lacinia urna, vel egestas nunc sapien ut nisl. Fusce placerat posuere est. Maecenas augue. In gravida leo vel massa. Aliquam eu nisi. Cras tincidunt aliquam arcu. Pellentesque habitant morbi tristique senectus et netus et malesuada fames ac turpis egestas. Aenean quis libero condimentum libero porttitor vestibulum. Pellentesque a nulla nec ipsum tristique molestie.

Curabitur sagittis quam vitae tortor. Nulla id pede. Praesent rutrum, ipsum quis porttitor pretium, nisl ipsum gravida est, commodo pulvinar augue nibh ac felis. Aliquam erat volutpat. Aenean ullamcorper feugiat dolor. Donec nulla mauris, elementum sed, varius ut, blandit vitae, purus. Aenean in nunc in justo sodales fringilla. Proin pulvinar. Maecenas ullamcorper. Vivamus dapibus consectetuer arcu. Nam dictum, erat non adipiscing pretium, lorem tellus pulvinar dui, ut condimentum massa est sit amet urna. Nulla facilisi.

Vivamus in dui. Morbi tristique nibh in purus. Nam erat. Curabitur euismod lacinia purus. Maecenas tempor libero quis dolor. Quisque gravida lorem vitae libero. Aliquam sodales, risus in tempor imperdiet, massa nisi auctor massa, vitae laoreet nunc erat in sapien. Etiam eu sapien nec enim consectetuer accumsan. Nunc non quam. Cras placerat justo quis nibh nonummy sollicitudin. Aenean vitae felis. Mauris id mi nec ligula ultricies vestibulum. Suspendisse luctus pretium justo. 

%\include{kapitel2}
%\include{kapitel3}
%\include{kapitel4}
%\include{kapitel5}
%\include{kapitel6}
%\include{kapitel7}
\cleardoublepage

% VERZEICHNISSE (Abbildungen, Tabellen)
% Literatur
\nocite{wiki:wissarbeit}
\bibliographystyle{alphadin}
\bibliography{Bachelorarbeit}
\cleardoublepage

% ERKLÄRUNG
\addcontentsline{toc}{chapter}{Formalities (DE)}

\chapter*{Eidesstattliche Versicherung}
\markboth{Formalities (DE)}{Formalities (DE)}
Hiermit versichere ich an Eides statt, dass ich die vorliegende Arbeit im
Bachelorstudiengang Software-System-Entwicklung selbstständig verfasst und
keine anderen als die angegebenen Hilfsmittel – insbesondere keine im
Quellenverzeichnis nicht benannten Internet-Quellen – benutzt habe. Alle Stellen,
die wörtlich oder sinngemäß aus Veröffentlichungen entnommen wurden, sind als
solche kenntlich gemacht. Ich versichere weiterhin, dass ich die Arbeit vorher nicht
in einem anderen Prüfungsverfahren eingereicht habe und die eingereichte
schriftliche Fassung der auf dem elektronischen Speichermedium entspricht.

\vspace{2cm}

\noindent Hamburg, den \uline{~~~~~~~~~~~~~~~~~~~~}~~~~~Unterschrift: \uline{~~~~~~~~~~~~~~~~~~~~~~~~~~~~~~~~~~~~~~~~~~~~~~~~~~}


\chapter*{Einstellung in die Fachbereichsbibliothek}
\markboth{Formalities (DE)}{Formalities (DE)}
Ich bin mit einer Einstellung dieser Arbeit in den Bestand der Bibliothek des Fachbereiches Informatik an der Uni Hamburg einverstanden.

\vspace{2cm}

\noindent Hamburg, den \uline{~~~~~~~~~~~~~~~~~~~~}~~~~~Unterschrift: \uline{~~~~~~~~~~~~~~~~~~~~~~~~~~~~~~~~~~~~~~~~~~~~~~~~~~}


\end{document}
